\begin{abstract}
This preprint addresses the frozen question Q4: whether there exist formal
No-Go theorems that rule out enterprise \emph{transformation} under explicitly
specified structural constraints, when enterprises are modeled as continua
$\KEA$ within the Ontology of Continua (\OC) and the executable specification
\kw{ETS.K\_EA.v1.1}. The contribution is strictly negative in scope.

We formulate impossibility results that arise from structural
incompatibilities among axes $\Ax$, admissible threshold configurations $\Th$,
boundary conditions $\dOm$, cycle-closure requirements $\Cyc$, and
phase-transition constraints governing the diagnostic states $\PHOK$, $\PHIN$,
$\PHCO$, and $\PHST$. Transformation is defined as a change of structural type
that cannot be represented as a trajectory entirely contained within the
original admissible state set $\Om(\KEA)$ under a fixed primitive signature.
This notion is formally distinguished from evolution, which denotes admissible
change within a preserved structural configuration.

Under this distinction, we show that for broad and explicitly characterized
classes of enterprise continua, transformation is formally impossible under
the given OC/ETS constraints, independently of time, effort, agency, or
intervention. The impossibility follows solely from the internal structure of
the continuum, rather than from contingent empirical or contextual factors.
For each No-Go result, precise falsifiability conditions are specified,
formulated entirely at the OC/ETS level, under which the corresponding
impossibility claim would fail.

The preprint is self-contained, non-prescriptive, and introduces no new
primitives, operators, or semantic reinterpretations beyond those already
defined in \OC\ and \kw{ETS.K\_EA.v1.1}.
\end{abstract}

\noindent\textbf{Keywords:}
enterprise continuum; structural constraints; No-Go theorems; impossibility
results; thresholds; phase logic; Ontology of Continua; $\KEA$
